% =====================================================================
% ONR_white_paper_5pp.tex -- AICHAINS white paper, ONR BAA N0001425SB001,
%                    Mathematical Data Science (Code 311)
%
% Submission-length version, Oct 29, 2026: cover page (unnumbered, not
% counted) + 5 body pages, the BAA limit. Derived from ONR_white_paper.tex
% (cover + 6 body pages) by compression; the argument and the PI's wording
% are preserved, the redundancy between the overview and the thrusts is not.
% Format: 8.5x11, 1-inch margins, single-spaced, Times New Roman 12 pt.
% Compile with: xelatex ONR_white_paper_5pp (twice, for cleveref).
%
% BAA-mandated layout: 1. Technical Concept (innovation + technical risk
% areas); 2. Future Naval Relevance; 3. Rough Order of Magnitude (ROM) Cost
% Estimate. The Operational Naval Concept is the primary use case in Section 2;
% the Operational Utility Assessment Plan is stated in risk (iv): seeded fault
% injection, retrospective industrial episodes, pre-registered criteria, and no
% field experiment (basic research).
%
% Figure: one panel, fully synthetic (redrawn from figs/panel_b_tikz_backup.tex):
% candidate incident chains in a hypothesis space. The NDA-gated
% agent-interface screenshot and the 1800x990 route-history rendering used in
% the 6-page version are not included here; at white-paper scale their labels
% are unreadable, and the screenshot needs Seagate's written clearance. Both
% remain in ONR_white_paper.tex for the PI to restore if desired.
%
% Placeholders to fill before submission: Co-PI email and exact title, SPS
% contact, submission date, UEI/CAGE, Co-PI effort in the ROM table.
% =====================================================================
\documentclass[12pt]{article}
% Strict font compliance: BAA requires Times New Roman 12 pt -- use the actual
% system Times New Roman via XeLaTeX/fontspec.
\usepackage[letterpaper,margin=1in]{geometry}
\usepackage{iftex}
\usepackage{url}
\urlstyle{same} % URLs in Times New Roman, not monospace
\ifXeTeX
  \usepackage{fontspec}
  \setmainfont{Times New Roman}  % real Times New Roman (BAA requirement)
\else
  % pdflatex fallback: metric-compatible Times clone; FINAL submission should
  % be compiled with xelatex so the embedded font is genuine Times New Roman.
  \usepackage{newtxtext,newtxmath}
\fi
\usepackage{graphicx}
\usepackage{tikz}
\usetikzlibrary{arrows.meta,positioning}
\usepackage[font=normalsize]{caption} % 12 pt captions, strict BAA reading
\usepackage{cleveref}
\usepackage{titlesec}

% Plain, compact section headings (12 pt bold, word-processor style)
\titleformat{\section}[block]{\normalsize\bfseries}{}{0em}{}
\titlespacing*{\section}{0pt}{6pt plus 1pt}{4pt}

\emergencystretch=1em
\raggedbottom

% Engine-portable delta-t label for the figure (XeLaTeX renders the real
% Times New Roman glyph; pdfLaTeX falls back to math mode).
\ifXeTeX
  \newcommand{\deltat}{\textit{Δt}}
\else
  \newcommand{\deltat}{$\Delta t$}
\fi

\begin{document}

\thispagestyle{empty} \begin{center} \textbf{WHITE PAPER SUBMITTED IN RESPONSE TO}\\[4pt]
\textbf{ONR BAA N0001425SB001 (as amended by Amendment 0008, Sep 2026)}\\[4pt]
Program: Mathematical Data Science (Code 311) --- Dr.~Predrag Neskovic \end{center}

\vspace{0.6in}

\noindent\textbf{Title:} AICHAINS: AI Causal Hypothesis Analysis and Inference System---Certified Causal Root-Cause Analysis for Complex Systems with Guarantees

\vspace{0.25in}

\noindent\textbf{Offeror:} Purdue University, West Lafayette, Indiana\newline \textbf{PI (technical contact):} Bruno Ribeiro, Associate Professor, Computer Science, Purdue University\newline {ribeirob@purdue.edu} \textbar{} \url{https://www.cs.purdue.edu/homes/ribeirob/} \newline \textbf{Co-PI:} Sanjay Rao, Professor, Computer Science, Purdue University\newline {[}email, exact title --- verify{]} \newline \textbf{Administrative contact:} Purdue SPS {[}name/phone/email --- request{]}\newline \textbf{Proposed period of performance:} 3 years, no options.

\vspace{0.25in}

\noindent Rough cost estimate: \$0.68M total over 3 years (about \$155K/year direct excluding student fee remissions and F\&A-excluded items, plus indirect at Purdue F\&A 57\% of MTDC; see Section 3).

\vspace{0.25in}

\noindent Topic areas addressed (per the Mathematical Data Science program page): \begin{itemize} \item Dynamical properties of networks and determination of causal effects \item Computational algorithms that scale to large networks with performance guarantees \end{itemize}

\vspace{0.25in}

\noindent Date: {[}submission date{]} \textbar{} Purdue UEI/CAGE: {[}request from SPS{]}

\newpage \setcounter{page}{1}

\section*{1. Technical Concept}

\begin{figure}[htbp]
\centering
% NDA-safe synthetic schematic (figs/panel_b_tikz_backup.tex, redrawn at
% natural size). The NDA-gated agent-interface screenshot of the 6-page
% version is not used here; swap it back only with Seagate's written clearance.
\begin{tikzpicture}[font=\small, line cap=round,
  evt/.style={draw=black!60, rounded corners=1.5pt, fill=blue!8, align=center, inner sep=2.4pt},
  tgt/.style={draw=black!70, rounded corners=1.5pt, fill=red!12, align=center, inner sep=2.8pt, thick},
  arr/.style={-{Latex[length=1.3mm]}, draw=black!60},
  dt/.style={font=\small\itshape, text=black!55},
  lab/.style={font=\small\bfseries, text=black!65}]
% AICHAINS example (synthetic; mirrors the deployed encoding at problem-class
% level: candidate incident chains with fuzzy alternative sets, by lead time)
\node[evt] (h3) at (0,1.90) {rework skipped (lots \{L2, L4\})};
\node[lab, left=1mm of h3] {H\textsubscript{3}};
\node[evt] (h2) at (0,1.00) {toolset \{M1, M3\} at polish};
\node[lab, left=1mm of h2] {H\textsubscript{2}};
\node[evt] (h1a) at (0,-0.30) {PM skipped at etch};
\node[lab, left=1mm of h1a] {H\textsubscript{1}};
\node[evt, right=3mm of h1a] (h1b) {chamber drift};
\node[tgt] (tg) at (6.6,0.90) {defect at final\\ inspection};
\draw[arr] (h1a) -- (h1b);
\draw[arr] (h1b.east) to[out=0,in=200] (tg.south west);
\draw[arr] (h2.east) to[out=0,in=185] node[dt, pos=0.55, above=0.6pt] {\deltat{} = 9 steps} (tg.west);
\draw[arr] (h3.east) to[out=0,in=155] node[dt, pos=0.55, above=0.4pt] {\deltat{} = 14 steps} (tg.north west);
\node[dt, align=center, below=2mm of h1a.south] {alternatives sorted by lead time;\\ incompatible survivors: abstain/escalate};
\end{tikzpicture}
\caption{Candidate incident chains in a hypothesis space, each a causal mechanism with events, edges, and timing; alternatives are sorted by lead time ($\deltat$ is distance to the target step), and incompatible survivors trigger abstention or escalation.}
\label{fig:aichains}
\end{figure}

\textbf{The problem.} When a complex system degrades---a manufacturing line, a logistics network, a ship's propulsion plant, a communication network---the operational question is which node, change, or propagation chain caused the incident, and what it will affect next. Answering it requires correlating sensor telemetry, logs, tickets, and change records over time and across organizational boundaries; performed manually, this correlation is slow. In a recent diagnosis at Seagate's semiconductor manufacturing plant, an engineering team spent weeks on a single recurring defect, rerouting production while defects continued, without reaching closure.

Our industry experience with RCA (Ford, Seagate, Amazon) has shown us that diagnosis suits agentic automation because it is iterative and hypothesis-driven: the analyst queries a data source, forms a suspicion, follows it into another source, and refines the hypothesis. The standard approach for such agents is retrieval-augmented generation (RAG): a vector index over \emph{text}---tickets, manuals, knowledge-base articles---queried at each RCA step. Current agentic and RAG architectures are mismatched to systems diagnosis in three respects. First, existing diagnostic tools are designed for human consumption: they collapse the space of candidate causes into a few point diagnoses, discarding the \emph{epistemic uncertainty}---uncertainty about what has not been observed---arising from uninstrumented steps, unlogged interventions, and dependencies absent from the recorded system map.

Second, multi-hop reasoning is difficult. Tickets omit technical context, change logs do not capture dependencies or all potential downstream consequences, and linking a symptom to a cause may traverse a chain of undocumented relationships. For example, a defect found at inspection may have been introduced three process steps earlier by a tool, on specific machines, and only for lots that skipped an intermediate rework step. Third, current RAG operates on text and on sources convertible to text, where documents and knowledge bases are written for retrieval. A plant's diagnostic evidence---sensor tensors, dependency graphs, event chains, and unstructured logs---is not as easily retrievable with existing methods. What the agent must retrieve is a structured \emph{candidate explanation for an issue}, together with related smaller issues, past or present, and the process changes that may precede it. Existing agentic RAG solves only about 11\% of failure cases on the OpenRCA benchmark,\footnote{J.~Xu, Q.~Zhang, Z.~Zhong, S.~He, C.~Zhang, Q.~Lin, D.~Pei, P.~He, D.~Zhang, and Q.~Zhang, ``OpenRCA: Can Large Language Models Locate the Root Cause of Software Failures?,'' in \emph{Proc.\ Int.\ Conf.\ on Learning Representations (ICLR)}, 2025.} and a tree-search diagnostic agent that exceeds 90\% accuracy on an open benchmark loses about 26 points on production incidents.\footnote{A.~Naakka, Y.~Wang, and M.~V. M\"{a}ntyl\"{a}, ``LATS-RCA: Language Agent Tree Search for Root Cause Analysis in Microservices,'' in \emph{Proc.\ 52nd Euromicro Conf.\ on Software Engineering and Advanced Applications (SEAA)}, 2026; arXiv:2605.03505.} These limitations have a common root: RCA is a \emph{causal inference} problem---genuine cause chains versus spurious correlation under partial observability---yet formal causal machinery and practical multimodal diagnosis pipelines have developed largely in isolation.



\noindent
\textbf{The approach.} A complex system is a network of interacting components. Current diagnostic tools, including our deployed hypothesis-space system, establish \emph{associations}: they narrow a population of candidate explanations against the evidence, but a surviving hypothesis is a dependence pattern. The open problem is the passage from associations in hypothesis spaces to identifying causation: a candidate causal reason may already be present in the hypothesis space, and the task of the proposed research is to identify which mechanism the data support.



We claim certified RCA for such systems is achievable if we (1) represent the state of diagnosis as an explicit \emph{hypothesis space}---a structured, editable population of candidate causal incident chains that preserves, rather than prematurely collapses, epistemic uncertainty; (2) define \emph{operators} over these objects---intersection, temporal association, corpus-scale matching---with performance guarantees; and (3) adjudicate the surviving hypotheses through \emph{causal verification}: causal discovery, from associational data, within the subpopulations (strata) that the hypothesis-space tree's splits define; interventional validation where feasible; and conformal abstention that bounds how often the system outputs a wrong or spurious cause.

Each step is a research contribution that, to our knowledge, no existing diagnostic system offers. (1) Diagnosis state is carried as a data-derived, editable population of candidate causal incident chains, which can also absorb hand-authored procedures (playbooks, dependency maps, correlation rules) as further hypotheses. (2) Operators match a hypothesis space to related hypothesis spaces under declared identifiability conditions. (3) The intersection of hypothesis spaces is an associational filter that retains the true cause whenever the true cause is among the candidates, and a separate causal layer determines which survivors are causal.

We will develop AICHAINS (AI Causal Hypothesis Analysis and Inference System) to implement this approach. It addresses the program's two challenge areas: determination of causal effects in networks, making the causal chain the primary object of inference; and algorithms that scale to large networks with performance guarantees, through auditable operators that carry distribution-free or structural guarantees.

\textbf{(1) Hypothesis spaces as populations of candidate causal chains.} A hypothesis space (\Cref{fig:aichains}) is the set of plausible failure modes consistent with observed data, encoded as time-indexed, text-annotated graphs satisfying three properties: (P1) capture relationships between hypotheses (refinements, subsumptions); (P2) stay compact---tree-like encodings over equivalent multi-feature splits avoid the combinatorial explosion of flat cross-products; (P3) preserve epistemic uncertainty---premature collapse of the space produces confident wrong diagnoses. Our preliminary deployed study illustrates (P3): an interpretable tree model over a manufacturing route-history dataset produced seven mutually incompatible explanations of one split, spanning hundreds of upstream operations. Reporting only the single best split would have yielded an incorrect root cause. AICHAINS advances this representation in a causal direction: each hypothesis is an explicit \emph{incident chain}, a candidate causal mechanism with events, propagation edges, and timing. We will develop techniques that find in the space the Markov equivalence classes supported by the data---those consistent with the observed conditional independence statements---and Thrust (3) then adjudicates the survivors. We will design two-part MDL encodings---a structure code over edges and timing, a residual code over events a chain leaves unexplained---so that model selection prunes the space while equivalent splits mark what the data cannot distinguish. Chain structure is elicited jointly from data and domain experts through three-valued edge feedback (certainly, certainly not, maybe), extending our group's human-in-the-loop causal discovery line. The contribution is the representation: an editable, uncertainty-preserving population of candidate causal chains that no current diagnostic system maintains.

\textbf{(2) RAG over hypothesis spaces.} Text RAG matches a text query to documents, images, or videos. AICHAINS RAG matches the current hypothesis space to related hypothesis spaces: an elevated fail rate---whether or not it triggered a response---can be matched to similar episodes, past or still unfolding, and to the process changes that preceded it. An operator is \emph{typed}: its specification declares the structure of the hypothesis spaces it consumes and returns, and carries an \emph{identifiability statement}---what the output licenses (dependence structure, temporal precedence) and what it cannot (causation, absent the verification of Thrust (3)). \emph{Intersection} retrieves the evidence for a hypothesis: it aligns symptom spaces (from live telemetry) against cause spaces (from change logs, tickets, association graphs) via graph transformers equivariant to the representation's symmetry group, fused with text encoders by cross-attention. We will characterize when the intersection is injective, so that every surviving hypothesis traces to exact supporting artifacts. \emph{Temporal association} retrieves a hypothesis's history: it composes static equivariant snapshot encodings through a graph-then-time framework into variable-size embeddings that localize when a hypothesis first appears, disappears, or generalizes. Temporal precedence is computed from the data. \emph{Corpus-scale matching} retrieves entire past hypothesis spaces: (set-of-matrices)/(set-of-matrices) matching over permutation-equivariant node-level embeddings, scored by subgraph-isomorphism networks, pre-trained label-free from LLM reasoning traces of resolved incidents via Gumbel-Sinkhorn alignment. The operators are associational; they become causal through Thrust (3), where the discovery assumptions (faithfulness, no unmeasured confounders) are stated explicitly and discovery runs within the strata that the hypothesis tree's splits define, keeping constraint-based and score-based methods tractable on high-dimensional network data. The contribution is the retrieval operators: matching structured hypothesis spaces to each other with guarantees.

\textbf{(3) Certified causal verification and abstention.} The hypothesis space is the agent's state, the operators are transitions over it, and the agent's memory is a corpus of prior hypothesis spaces retrieved by corpus-scale matching. Every accepted diagnosis is accompanied by a machine-checkable chain: each surviving incident chain must (a) trace to provenance-verified artifacts, (b) pass structural invariants (typing, dataflow consistency, strictly-decreasing consistency residuals---discharged with an SMT solver), (c) pass temporal precedence against the incident onset, and (d) survive \emph{interventional validation} where feasible---seeded fault injection in test systems and natural experiments in operation. Reliability scores are mapped to decisions through conformal prediction with ordering-invariant score functions, yielding distribution-free coverage and a calibrated \emph{abstain-or-escalate} verdict: when the surviving space still contains incompatible chains with incompatible interventions, the system escalates to a human with the ambiguity made explicit rather than an unsupported verdict. Faithfulness probes (counterfactual ablation of claimed causal edges) guard against post-hoc rationalization. The contribution is the causal layer: deciding which chains surviving the intersection are causal, with bounded error or an explicit abstention. \Cref{fig:aichains} shows a hypothesis space of candidate incident chains.

\textbf{Feasibility.} The proposed research builds on validated prior work: the PI's group authored the underlying representation theory (Holographic, ICLR 2025; STAGE, ICML 2025; All-IN, ICLR 2026), the temporal equivariance framework behind temporal association, a human-in-the-loop causal discovery program, and PRISM, a video streaming RCA system the PI's and Co-PI's groups built with a major video-streaming provider. A system implementing the tree-based hypothesis-space diagnosis loop---the associational stage of the proposed program---is deployed in Ford's INSIGHT tool under a previous grant, and a version for semiconductor manufacturing is in development with Seagate (data under NDA). The LLM agent only selects and parameterizes deterministic analysis tools, accessed through typed tool interfaces with read-only enforcement; every number an engineer sees comes from the analysis code rather than from the model. The trees use positive set-based splits, report the exhaustive set of equivalent splits at every node, and, through leakage gates, drop any candidate event that occurs after the target for even one unit (which is a consequence rather than a root cause). Confidence is scored as an entropy capture ratio against a gradient-boosting ceiling with a Nadeau--Bengio-corrected significance test. Results are deterministic: a cache keyed on data, parameters, and code digests returns identical answers in 0.02 s versus about 25 s of fresh analysis at production scale. The system is air-gapped with a local LLM. It was built to four design requirements (provenance, adaptability, reliability, uncertainty quantification) and evaluated against engineering-team investigations and synthetic data with by-construction root causes. PRISM found multiple simultaneous issues, including tail incidents missed by strong baselines, on real production snapshots.\footnote{C.~Bothra, X.~Jiang, Y.~Zhang, Z.~Akhtar, S.~G.~Rao, and B.~Ribeiro, ``PRISM: Localizing performance anomalies for large scale video services,'' manuscript in preparation, 2026.}

\textbf{Technical risk areas and mitigations} (per the BAA). (i) \emph{Discovery under confounding:} constraint-based discovery on observational fab or network data may be unfaithful. Mitigation: discovery runs within the hypothesis tree's strata, under stated assumptions, is validated interventionally where feasible, and ranks candidates rather than certifying them; it remains subordinate to provenance and precedence checks. (ii) \emph{Hypothesis explosion:} intersecting multimodal spaces may explode combinatorially. Mitigation: compactness-by-construction encodings with tunable resolution let the agent control search breadth under explicit query budgets. (iii) \emph{Elicitation noise:} expert or LLM edge feedback is fallible. Mitigation: three-valued feedback feeds a contradiction registry; edges carry confidence and provenance; the conformal layer absorbs residual uncertainty. (iv) \emph{Integration risk:} the proposed operators and the causal-verification layer are yet to be developed, and their integration with the deployed associational stage is unproven; the seeded fault-injection protocol, retrospective industrial episodes, and pre-registered success criteria make the claim falsifiable.

\section*{2. Future Naval Relevance}

Root-cause analysis under time pressure is a recurring naval problem. A frigate's propulsion plant, a shipyard's build schedule, a semiconductor fab's yield line, and a sustainment network are graphs of interacting components whose failures must be explained from sparse, noisy, multimodal evidence, typically without ground truth for the most consequential failure modes. AICHAINS is designed for this setting; we develop it across naval platform health, shipyard and microelectronics manufacturing, and fleet logistics.

\textbf{Naval platform health and condition-based maintenance (primary use case).} A ship's propulsion plant is a dependency graph of sensors, actuators, and subsystems; a combat system, a chain of processing, power, and cooling dependencies. When a casualty occurs, the questions are which chain of events produced the anomaly, which subsystems it will affect, and what to inspect first. AICHAINS maintains and edits populations of candidate failure-propagation chains over the platform's dependency graph. Intersection aligns live telemetry against maintenance and configuration changes. Temporal association localizes when a propagating degradation began. The certified chain gives the watch team a defensible answer with a bounded error rate or, when evidence is insufficient, an explicit statement of the surviving ambiguity.

\textbf{Naval manufacturing, shipyard repair, and trusted microelectronics (bridge domain).} The PI's group operates a deployed, regression-tested RCA system that implements the proposed hypothesis-space loop in Ford's INSIGHT tool (from a previous grant), and a version for semiconductor manufacturing is in development with Seagate (data under NDA; described at problem-class level). Shipyard workflows, depot overhaul, naval additive manufacturing, and trusted microelectronics production are process-dependency graphs: defects propagate across stages, tools, and lots, and yield is lost to failures invisible in the final product. Transfer to the naval setting is supported by attribute-domain-agnostic representations rather than by domain-specific engineering. This proposal adds the formal causal machinery that moves these systems from association to causation.

\textbf{Fleet sustainment and logistics networks.} A strike group's sustainment is a multi-echelon supply network; parts flow and maintenance scheduling form a transport network over the platform graphs. Delay, disruption, and root-cause prediction reduce to the same problem: localizing fault propagation on a process-dependency graph from noisy, label-scarce observations, with the same certified-trace and abstention guarantees.

\textbf{Program alignment.} These use cases fall within the program's Complex Networks concentration area (basic research on the dynamical properties of networks and the determination of causal effects, and on algorithms that scale with performance guarantees), and they bear on multi-modal information integration and decision making under uncertainty. The research is basic: whether a population of associational hypotheses can be certified as causal---identifiability, compactness, and the guarantee theory required for trustworthy agentic diagnosis---is a foundational question for any naval deployment of diagnostic AI.

\section*{3. Rough Order of Magnitude (ROM) Cost Estimate}

Faculty effort: PI (0.46 summer person-months/year, 3.85\% effort), Co-PI Sanjay Rao (effort to be set --- [verify]; current figures exclude Co-PI salary), two PhD students (4.5 academic-year months plus 3 summer months in Years 1--2; academic year only in Year 3), no postdoc. Direct costs: salaries and fringe benefits, graduate fee remissions, HPC compute (Purdue Gautschi Center), cloud storage, travel (1--2 conferences and one ONR review meeting/year). Indirect costs at Purdue's negotiated rate of 57\% of MTDC (DHHS-approved 10/03/2022). A full R\&R budget follows in any invited application.

\begin{center} \normalsize \renewcommand{\arraystretch}{0.92} \begin{tabular}{lccc} \hline Year & Direct & Indirect (57\% MTDC) & Total \\
\hline
Y1 & \$168,560 & \$77,041 & \$245,601 \\
Y2 & \$173,566 & \$79,895 & \$253,461 \\
Y3 & \$122,869 & \$54,041 & \$176,910 \\
Total & \$464,995 & \$210,977 & \$675,972 \\
\hline \end{tabular} \end{center}

\end{document}
